\par\indent\SourceLeaf{214}%
Insofern ein Dreiecksnetz sich doch immer nur über einen sehr mässigen Theil der Erdoberfläche erstreckt, lässt sich der erwähnte Zweck noch vollkommner erreichen, wenn man die allgemeine Auflösung noch etwas generalisirt, und nicht \(f\upsilon=\upsilon\), sondern \(f\upsilon=\upsilon+\text{Const.}\) annimmt. Offenbar würde hiedurch gar nichts gewonnen, wenn man dieser Constante einen reellen Werth beilegte, weil dadurch lediglich \(T\) und \(t\) um diese Constante verschieden, also nur die Anfangspunkte der Längen ungleich werden würden. Allein ganz anders verhält es sich, wenn man der Constante einen imaginären Werth beilegt. Setzt man dieselbe \(=-i\log k\), so wird
\[
T=t,\qquad \mathrm{tang}\tfrac{1}{2}U=k\,\mathrm{tang}\tfrac{1}{2}w\mathord{.}(\frac{1+\varepsilon\cos w}{1-\varepsilon\cos w})^{\frac{1}{2}\varepsilon}
\]
\par\indent Um hier über den zweckmässigsten Werth von \(k\) entscheiden zu können, müssen wir vor allen Dingen das Vergrösserungsverhältniss bestimmen.

Es wird hier, in den Zeichen des 5. und 7. Artikels
\[
\begin{aligned}
n&=aa\sin u^2\\
N&=AA\sin U^2\\
\varphi\upsilon&=1
\end{aligned}
\]
\par\indent Also
\[
\begin{aligned}
m&=\frac{A\sin U}{a\sin u}=\frac{A\sin U}{a\sin w}\mathord{.}\BareRoot{(1-\varepsilon\varepsilon\cos w^2)}\\
&=\frac{A}{a}\mathord{.}\frac{k(1-\varepsilon\varepsilon\cos w^2)^{\frac{1}{2}+\frac{1}{2}\varepsilon}}{\cos\frac{1}{2}w^2(1-\varepsilon\cos w)^{\varepsilon}+kk\sin\frac{1}{2}w^2(1+\varepsilon\cos w)^{\varepsilon}}
\end{aligned}
\]
\noindent welches Verhältniss also bloss von der Breite abhängt. Die möglichst geringste Abweichung von vollkommner Aehnlichkeit erhält man, wenn man \(k\) so bestimmt, dass \(m\) für die äussersten Breiten gleich grosse Werthe erhält, wodurch von selbst \(m\) bei der mittlern Breite seinem grössten oder kleinsten Werthe sehr nahe sein wird. Bezeichnet man die äussersten Werthe von \(w\) durch \(w^0\) und \(w'\), so erhält man auf diese Weise
\[
k=\BareRoot{\frac{\displaystyle\frac{\cos\frac{1}{2}{w^0}^{2}(1-\varepsilon\cos w^0)^{\varepsilon}}{(1-\varepsilon\varepsilon\cos {w^0}^{2})^{\frac{1}{2}+\frac{1}{2}\varepsilon}}-\frac{\cos\frac{1}{2}{w'}^{2}(1-\varepsilon\cos w')^{\varepsilon}}{(1-\varepsilon\varepsilon\cos {w'}^{2})^{\frac{1}{2}+\frac{1}{2}\varepsilon}}}{\displaystyle\frac{\sin\frac{1}{2}{w'}^{2}(1+\varepsilon\cos w')^{\varepsilon}}{(1-\varepsilon\varepsilon\cos {w'}^{2})^{\frac{1}{2}+\frac{1}{2}\varepsilon}}-\frac{\sin\frac{1}{2}{w^0}^{2}(1+\varepsilon\cos w^0)^{\varepsilon}}{(1-\varepsilon\varepsilon\cos {w^0}^{2})^{\frac{1}{2}+\frac{1}{2}\varepsilon}}}}
\]
\noindent Um zu erfahren, bei welcher Breite \(m\) seinen grössten oder kleinsten Werth erhält, haben wir
